% BHU PYQ -- Manually transcribed & error-corrected
\documentclass[11pt,a4paper]{article}

\usepackage[a4paper,top=2.5cm,bottom=2.5cm,left=2.8cm,right=2.8cm,headheight=14pt]{geometry}
\usepackage{amsmath,amssymb}
\usepackage[T1]{fontenc}
\usepackage[utf8]{inputenc}
\usepackage{lmodern}
\usepackage{microtype}
\usepackage{fancyhdr}
\usepackage{enumitem}
\usepackage{hyperref}
\usepackage{lastpage}
\usepackage{parskip}

\hypersetup{colorlinks=true,urlcolor=black,linkcolor=black,
  pdftitle={BPT-101: Mechanics and Relativity -- BHU PYQ},pdfauthor={Banaras Hindu University}}

\pagestyle{fancy}\fancyhf{}
\renewcommand{\headrulewidth}{0.4pt}\renewcommand{\footrulewidth}{0.4pt}
\fancyhead[L]{\small   \textit{Banaras Hindu University}}
\fancyhead[R]{\small   \textit{BPT-101: Mechanics and Relativity}}
\fancyfoot[C]{\small Page  \thepage\ of \pageref{LastPage}}


\newlist{parts}{enumerate}{2}
\setlist[parts,1]{label=(\alph*),leftmargin=*,itemsep=5pt,topsep=3pt}
\setlist[parts,2]{label=(\roman*),leftmargin=*,itemsep=3pt,topsep=2pt}

\newcommand{\pts}[1]{\hfill   \textit{[#1]}}


\begin{document}


\begin{center}
  {\large\bfseries BANARAS HINDU UNIVERSITY}\\[2pt]
  {\normalsize B.Sc. (Hons.) Semester I Examination, 2023-24}\\[6pt]
  \rule{0.8\linewidth}{0.5pt}\\[6pt]
  {\large\bfseries Paper No. BPT-101: Mechanics and Relativity}\\[4pt]
  {\normalsize Time: 3 Hours\hfill Maximum Marks: 70}
\end{center}
\rule{\linewidth}{0.4pt}
\smallskip



\noindent   \textit{Instructions: Answer all questions. Symbols have their usual meanings.}
\bigskip




\noindent   \textbf{Question 1.} Answer any   \textbf{four} of the following questions: \pts{4 $ \times$ 3.5 = 14}

\begin{parts}
  \item According to an observer on Earth, a car covers $1\, \text{mile}$ in $50\, \text{s}$ on a straight stretch of highway. How long will it take to cover this distance according to an observer with his own clock on a space ship receding (moving away) from Earth with a speed of $0.95c$: (i) perpendicular to the line of motion of the car, and (ii) along the same line of motion of the car?
  \item A frame of reference $S'$ rotates with uniform angular velocity $\vec{\omega}$ with respect to a stationary frame $S$ having a common origin. Show that the relationship between the time derivatives is:
    \[ \left(\frac{d\vec{A}}{dt}\right)_S = \left(\frac{d\vec{A}}{dt}\right)_{S'} + \vec{\omega}  \times \vec{A} \]
    Obtain expressions for the centrifugal force and the Coriolis force acting on a particle moving in the rotating frame.
  \item A particle moves under the influence of a force $\vec{F}$ and has instantaneous velocity $\vec{v}$. Show that $\vec{F} \cdot \vec{v} = \frac{dT}{dt}$, where $T$ is the kinetic energy of the particle. ``Angular momentum and angular velocity of a rigid body are not always parallel.'' Justify this statement.
  \item Find the greatest length of a steel wire that can hang vertically without breaking. Breaking stress for steel is $7.9  \times 10^9\, \text{dynes/cm}^2$, and the density of steel is $7.9\, \text{g/cm}^3$. (Take $g = 980\, \text{cm/s}^2$).
  \item Obtain the theoretical limiting values of Poisson's ratio $\sigma$ using the relations between Young's modulus ($Y$), bulk modulus ($K$), and shear modulus ($\eta$).
\end{parts}

\medskip\rule{\linewidth}{0.2pt}




\noindent   \textbf{Question 2.}

\begin{parts}
  \item What do you understand by invariant quantities in the special theory of relativity? Derive the invariants for space-time interval ($ds^2$) and energy-momentum ($E^2 - p^2c^2$). \pts{7}
  \item Describe the Michelson-Morley experiment. Why is it called the most famous ``failed experiment''? \pts{7}
\end{parts}

\medskip\rule{\linewidth}{0.2pt}




\noindent   \textbf{Question 3.}

\begin{parts}
  \item A uniform circular disc of radius $R$ oscillates in a vertical plane about a horizontal axis. Find the distance of the axis of rotation from the center of the disc for which the period of oscillation is minimum. What is the value of this minimum period? \pts{7}
  \item A capacitor charged to $2\, \text{V}$ is discharged through a ballistic galvanometer, yielding a corrected deflection of $9.6\, \text{cm}$. If the current sensitivity is $4.54  \times 10^7\, \text{div/A}$ and the periodic time is $12\, \text{s}$, calculate the capacity of the capacitor. \pts{7}
\end{parts}
\hfill   \textbf{OR}

\begin{parts}
  \item Derive the Lorentz velocity transformation equations between two inertial frames of reference moving with relative velocity $V$ along the $x$-axis. \pts{7}
  \item Show from Lorentz transformations that two events simultaneous ($t_1 = t_2$) at different positions ($x_1 
e x_2$) in a reference frame $S$ are not necessarily simultaneous in another reference frame $S'$. \pts{7}
\end{parts}

\medskip\rule{\linewidth}{0.2pt}




\noindent   \textbf{Question 4.}

\begin{parts}
  \item If two capillaries of radii $r_1$ and $r_2$ and lengths $l_1$ and $l_2$ are joined in series, derive an expression for the rate of flow of liquid through the arrangement using Poiseuille's formula. \pts{8}
  \item Let $P$ be the stress (pressure) applied normally over an area $A$ of a body of volume $V$ such that its volume decreases by $v$. Prove that the work done per unit volume is $\frac{1}{2}  \times  \text{stress}  \times  \text{strain}$. \pts{6}
\end{parts}
\hfill   \textbf{OR}

\begin{parts}
  \item If a body falls freely through a height $h$ from rest, show that the horizontal displacement of the body due to the Coriolis force at latitude $\lambda$ is equal to:
    \[ x = \frac{2}{3} \omega \cos\lambda \sqrt{\frac{2h^3}{g}} \]
    where $\omega$ is the angular velocity of the Earth. \pts{8}
  \item What are elastic and inelastic collisions? Two clay balls collide head-on in a perfectly inelastic collision. The first ball has a mass of $0.5\, \text{kg}$ and an initial velocity of $4\, \text{m/s}$ to the right. The second ball has a mass of $0.25\, \text{kg}$ and an initial velocity of $3\, \text{m/s}$ to the left. What is the decrease in kinetic energy during the collision? \pts{6}
\end{parts}

\medskip\rule{\linewidth}{0.2pt}




\noindent   \textbf{Question 5.}

\begin{parts}
  \item Derive an expression for the twisting couple per unit angular twist for a hollow cylinder and explain why a hollow cylinder is stronger than a solid cylinder of the same length, mass, and material. \pts{8}
  \item Compare the loads required to produce equal depression in two beams made of the same material and having the same length and weight, with the only difference being that one has a circular cross-section while the other has a square cross-section. \pts{6}
\end{parts}

\vfill
\rule{\linewidth}{0.4pt}

\begin{center}\small
  \href{https://bhu-exam.vercel.app/}{Click here to visit our website}\\[4pt]
  \href{https://me-aryan.vercel.app/}{Click here to visit our contributor}\\[4pt]
  \href{https://quashan-me.vercel.app/}{Click here to visit our contributor}
\end{center}

\end{document}
